\newcommand{\APage}[3]{\clearpage\SetRegion{附录 A・指标体系与数据来源}{teal}{compass}\phantomsection\pdfbookmark[1]{#1：#2}{#3}\kicker{APPENDIX A / #1}\PageTitle{#2}{指标体系与数据来源}}
\newenvironment{ATable}[1]{\par\vspace{3pt}\begingroup\sffamily\fontsize{9.5}{14}\selectfont\setlength{\tabcolsep}{5pt}\noindent\tabularx{\linewidth}{#1}\toprule}{\bottomrule\endtabularx\endgroup\par\vspace{4pt}}
\APage{A.1 框架与行为}{从发展目标到可观察的协作}{app-a1}
\lead{OSDGs确定价值方向，OSDI组织观察指标；两者共同连接问题、证据与行动。}
开源发展指标体系（Open Source Development Indicators，OSDI）以开源发展目标（Open Source Development Goals，OSDGs）为导向，观察人才参与、技术供给、网络影响、国际协作与治理条件，为第二至第五章提供共同的概念和数据说明。

体系参照《信息技术 开源治理》系列标准及已有研究，结合OpenRank、ModelRank与人工智能时代的资产特征展开。\textbf{四个目标、十个观察方向与29项候选指标是不同层级；候选目录不表示所有指标均已完成本期测量。}正文使用的具体数据、窗口和限制以各图说明为准。

\topic{三项方法原则}
\textbf{科学性：}从行为与协作关系出发，区分规模、结构和影响。\par
\textbf{可操作性：}保留公开数据、算法与版本入口，说明可获取和可复算的范围。\par
\textbf{可比性：}先对齐对象、行为、窗口与汇总规则，再进行跨期、跨地区和跨主体比较。

\topic{五类核心协作行为}
\begin{ATable}{p{11mm}p{27mm}p{39mm}X}
缩写 & 中文名称 & 英文名称 & 行为含义 \\\midrule
OI & 创建议题 & Open Issue & 提出问题、功能建议或技术讨论。\\
ID & 议题讨论 & Issue Discussion & 在Issue或Pull Request中参与讨论和评论。\\
OPR & 提交代码请求 & Open Pull Request & 提交代码、文档等修改请求。\\
PRR & 代码评审 & Pull Request Review & 对修改请求进行审查、评论、批准或提出修改意见。\\
MPR & 合并代码请求 & Merge Pull Request & 修改请求经评审后正式合并至项目。\\
\end{ATable}
\src{表 A-1　核心协作行为。依据本报告指标说明与OpenDigger方法资料整理。平台事件映射和具体计权以相应版本为准。}
五类行为覆盖提出问题、讨论、提交、评审与合并。Star、Fork、点赞和下载另作信号，不能替代已发生的协作工作量。
\vfill\ReportNote{本附录集中解释指标与口径。候选指标目录见1.3；逐图来源与数据状态随正文保留，详细输入与核查记录随配套研究材料保存。}

\APage{A.2 参与与影响}{先区分活跃，再理解贡献与影响}{app-a2}
\topic{活跃参与：观察当期是否发生协作}
\begin{ATable}{p{36mm}X}
指标 & 定义与使用 \\\midrule
活跃开发者 & 统计期内至少发生一次纳入规则的有效开源协作行为的开发者账号。\\
活跃开发者数量 & 在明确平台与身份规则下，对满足活跃条件的账号去重计数。\\
活跃开源仓库 & 统计期内发生有效协作行为的仓库；若来源采用其他规则，须另作说明。\\
活跃开源仓库数量 & 满足对应规则的去重仓库数量，不直接等于项目数量。\\
\end{ATable}
\src{表 A-2　活跃参与指标。OpenRank含历史继承，得分大于零不能自动证明当期活跃。月度去重人数不能直接相加为年度人数。}
OpenDigger的Activity是对协作行为计权后形成的强度指标，与活跃人数、事件总量不同。原表的来源读数、清洗方式及聚合规则应分别保留，不能用显示层四舍五入代替方法核验。

\topic{OpenRank贡献度：在项目或社区内理解协作}
社区OpenRank基于项目内部的协作网络，综合行为及其连接关系，观察开发者的相对贡献。它可用于识别核心参与者，分析贡献在个体、群体与地区间的分布，并观察跨地域贡献流向。\textbf{网络贡献信号不等于对代码质量、工时或个人能力的直接测量。}

\topic{OpenRank影响力：在观测网络中理解位置}
全域OpenRank观察开发者与仓库在所覆盖协作网络中的位置。进一步讨论企业、基金会、地区或技术生态时，需要明确主体映射与聚合方式。
\begin{ATable}{p{31mm}X}
评价对象 & 主要观察含义 \\\midrule
开发者 & 跨项目参与、核心贡献与生态连接位置。\\
仓库或项目 & 协作网络中的连接与影响；项目需另行界定多仓边界。\\
企业或组织 & 关联资源与参与者形成的网络影响，不等于员工贡献。\\
开源基金会 & 关联或托管项目的网络表现，不直接代表治理成效。\\
国家或地区 & 按指定归属规则汇总的网络位置，不代表综合国力。\\
技术生态 & 人工智能、云原生、操作系统等样本的协作结构。\\
\end{ATable}
\src{表 A-3　OpenRank影响力的对象层级。不能仅因使用同一指标，就互换贡献来源、项目归属和组织控制权。}

\APage{A.3 OpenRank象限图}{两种视角，两个观察对象}{app-a3}
\lead{社区贡献度与全域影响力，分别回答不同层次的协作问题。}
\plate{figures/new/openrank_quadrants.pdf}{图 A-1　OpenRank的观察对象与汇总视角}{依据提供的OpenRank象限图重绘。轴用于区分概念，不表示连续量值；连线表示观察或汇总关系，不表示因果或具体计算公式。}
\textbf{先定网络，再定对象。}社区视角关注项目内部贡献；全域视角关注已覆盖协作网络中的位置。“全域”不表示已经覆盖全球全部开源活动。

\textbf{汇总须保留映射。}个体、组织、地区与技术生态的统计，都需固定对象集合与有效期，处理多仓、镜像和交叉归属；不能默认把仓库分数简单相加就得到项目分数。

\APage{A.4 模型与人才}{关系、席位与人才代理信号}{app-a4}
\topic{ModelRank：观察模型派生与复用网络}
ModelRank用于观察模型在派生与复用网络中的影响。一个模型的网络位置既与复用关系有关，也与下游节点的连接有关。本报告沿用来源月榜；算法版本、输入边构造与归一化尚未完整独立重建，不将未经核对的计算细节写成统一公式。
\begin{ATable}{p{35mm}X}
指标 & 定义与使用 \\\midrule
ModelRank模型影响力 & 来源评价网络中的模型影响信号，不等于模型能力测试或实际使用效果。\\
ModelRank排名 & 同一月份、同一榜单范围内，按来源ModelRank排列的相对位置。\\
头部模型数量 & 特定主体进入指定Top K范围的模型版本数；同一家族可有多个版本。\\
\end{ATable}
\src{表 A-4　模型影响指标。3.6使用2026年6月Top 100月榜；5.2区分主期末与期后对照。Top 20、Top 100与Top 1000相互包含，不能相加。}
模型卡声明的基础模型关系，是复用研究的入口，并不意味着训练过程已被独立核验。下载、点赞与网络影响分别描述不同信号，不互换为用户数、部署量或公共成效。

\topic{开源人才与贡献多样性}
\begin{ATable}{p{35mm}X}
指标 & 定义与使用 \\\midrule
开源人才 & 沿用来源前5\%规则识别的高OpenRank开发者群体，作为核心人才的代理信号。\\
开源人才数量 & 在指定总体、时期和阈值下，归属于某地区或城市的入选账号数。\\
国际贡献占比 & 指定项目集合获得的境外开发者贡献，占全部已纳入贡献的比例。\\
贡献来源多样性 & 贡献在国家、地区或组织间的分布；使用具体比例或结构指标时分别说明。\\
\end{ATable}
\src{表 A-5　人才与多样性指标。人才筛选总体、累计窗口和并列处理尚需与来源规则逐项对应；不把前5\%解释为综合能力认证。}
贡献流向同时涉及开发者所在地与项目归属。二者不是同一地理概念；来源多样、跨境比例较高，也不能单独证明治理更加包容或合作效果更好。
\vfill\keypoint{比较前先核对}{筛选总体、对象身份、统计月份、未知项与截榜规则。}

\APage{A.5 治理与数据}{明确资产类型，也明确资料来源}{app-a5}
\begin{ATable}{p{42mm}X}
治理与数字资产指标 & 定义与使用 \\\midrule
基金会活跃仓库数量 & 在指定名录与映射下，由基金会关联或托管且满足活跃规则的仓库数。\\
基金会项目影响力占比 & 在相同网络、时期与聚合规则下，基金会关联项目分值占明确总体的比例。\\
许可证采用量 & 按可识别许可证统计的活跃仓库数；无记录与未采用分开处理。\\
有效活跃仓库数 & 统计期内具有实际协作行为、且满足数据清洗规则的代码仓库数。\\
开源模型数量 & 指定平台与快照中的公开模型记录；公开可见不自动等于完整开源。\\
模型发布单位数量 & 发布公开模型的账号或组织数；账号与法律主体不等同。\\
开源数据集数量 & 指定平台与范围内的公开数据集记录；开放条件另行核对。\\
数据集发布单位数量 & 发布公开数据集的账号或组织数，保留身份映射范围。\\
\end{ATable}
\src{表 A-6　治理与数字资产指标。代码、模型和数据集分别计量；基金会支持关系、正式托管与生态关联分别识别。}

\topic{主要数据来源与访问层}
\begin{ATable}{p{27mm}p{48mm}X}
资料类型 & 主要来源 & 用途 \\\midrule
开源协作行为 & OpenDigger及各平台公开记录 & 活跃、贡献与网络影响观察。\\
国内平台补充 & Gitee、AtomGit等 & 对应覆盖下的开发者、仓库和项目观察。\\
模型与数据集 & Hugging Face & 资源、发布账号、声明关系及模型榜单。\\
基金会资料 & 公开名录及关联仓库 & 主体边界、项目组合与治理关系。\\
许可证资料 & 仓库元数据、SPDX标识及相关记录 & 采用结构与开放条件核对。\\
地理资料 & 公开所在地字段及标准化地名库 & 国家、省市等地域映射。\\
\end{ATable}
\src{表 A-7　主要资料来源。实际平台覆盖和提取版本随各图说明。}
\textbf{OpenDigger是数据和指标基础，OpenShare是访问和应用这些数据的产品层。}两者不是两个独立证据来源。OpenWorld的GOSD专用入口采用报告期，普通入口采用自然年；取数时须明确入口和对象清单。

\APage{A.6 时间与对象}{让每一次比较保持同一口径}{app-a6}
\keypoint{默认主统计期}{2025-07-01 至 2026-06-30。其他窗口、历史材料和期后观察均在对应图文处另行标注。}
\topic{报告年度、自然年度与快照分别阅读}
报告年度不随外部资料的发布日期改变。引用GitHub年度窗口、Hugging Face夏季观察或其他研究时，保留来源的实际起止时间；快照按提取日标注，月榜按具体月份标注。仅写“截至2026年7月”，不足以证明覆盖主统计期。

\topic{先在平台内去重，再核对跨平台身份}
开发者以“平台＋稳定用户ID”、仓库以“平台＋稳定仓库ID”为基本识别方式。跨平台合并须有经核验的身份或镜像映射，不能仅凭同名去重。机器人与服务账号的识别规则、名单版本和排除范围应留存，不把局部规则表述为全平台已完成核验。

\topic{项目标签可能包含多个仓库}
固定项目—仓库清单、版本与生效时间，分别记录转移、拆分、镜像和交叉归属。项目、仓库、组织及基金会可有重叠；汇总前须消除重复计算，不能将不同榜单或样本直接相加。

\topic{未知地域不等于没有参与}
开发者所在地来自公开资料的标准化映射，不等于国籍。项目依据指定的发起方或主要维护组织规则归属，也不等于贡献来自该地。未知项保留在适用的全球总体中并单列；地区细分另说明可识别覆盖。历史分析优先使用相应历史快照，不能用当前资料自动还原历史位置。

\topic{平台覆盖与核查状态如实说明}
Gitee、AtomGit等补充仅在有依据的对应范围纳入；缺少可靠省市资料时不进入城市级分析。缺月、路径失效、对象差异与数值差异分别处理，不补零、不把局部一致扩大为整表认证。

第四章4.3、4.5、4.6的247条固定记录中，116条同时满足身份、12个月数据和原值一致；其余仍有缺口或差异。4.6是PyTorch生态观察样本，不等于基金会正式托管名单。
\vfill\ReportNote{资料依据：本报告宪法、获批主张与核查记录，OpenDigger方法归档，以及各章固定输入。证据状态不会因视觉整理或全书合并而改变。}

\APage{A.7 术语对照}{共同概念，保留各自的含义}{app-a7}
\begin{ATable}{p{59mm}X}
中文名称 & 英文名称 \\\midrule
开源发展目标 & Open Source Development Goals（OSDGs）\\
开源发展指标体系 & Open Source Development Indicators（OSDI）\\
活跃开发者 & Active Developer\\
活跃开源仓库 & Active Open Source Repository\\
OpenRank贡献度 & OpenRank Contribution Impact\\
OpenRank影响力 & OpenRank Influence\\
社区OpenRank算法 & Community OpenRank Algorithm\\
全域OpenRank算法 & Global OpenRank Algorithm\\
开源人才 & Open Source Talent\\
国际贡献占比 & International Contribution Share\\
开源数字资产 & Open Source Digital Asset\\
ModelRank模型影响力 & ModelRank Model Influence\\
模型派生与复用网络 & Model Derivation and Reuse Network\\
近30天下载量 & Downloads in the Last 30 Days\\
累计点赞量 & Cumulative Likes\\
\end{ATable}
\src{表 A-8　中英文术语对照。沿用报告中的表达，不据此改变上游字段、算法版本或统计定义。}
\vspace{12pt}\RunTitle{从指标返回问题}
人数回答谁参与，行为描述如何协作，网络指标刻画连接位置，资产记录呈现可见供给。治理与公共价值还需要具体版本、实际采用、维护责任、持续成本和结果证据。

这些概念可以互相补充，但不能互相替代。\textbf{可追溯的比较，始于明确的对象、时间和来源。}
\vfill\noindent\includegraphics[width=\linewidth]{figures/new/appendix_dimensions.pdf}
\src{四个目标、十个观察方向，是全书共同的阅读坐标；不是算法权重或证据等级。}
